\subsection{The dynamic model, and its equilibrium reduction}
\label{sec:dynamic}
%==========================================================================

\emph{One model, at the resolution where it is a differential equation. Everything after
this section is a reduction of it, and the reductions are named for what they do rather
than lettered.}
%==========================================================================

\subsubsection{What the model contains}
\label{sec:model_C_added}

The state is the plasma marker, the enzymatic channel is Michaelis-Menten in it, the therapy
enters as a time-varying increment to the capacity driven by a known dosing schedule, and
the propiogenic load is an unknown input:
\begin{equation}
\dot m \;=\; L(t) + L_0 \;-\; \bigl(\theta + \Delta(t)\bigr) g(m) \;-\; \alpha(m),
\qquad g(m) = \frac{m}{K_m+m}.
\label{eq:parent_orig}
\end{equation}
The rest of this section is about $\alpha$, the alternative clearance channel, because that
is the only term the three candidate structures ever disagreed about. Two facts about
\eqref{eq:parent} are settled first and are not revisited: $\theta$ absorbs the maximal
velocity, so it \emph{is} the residual $V_{\max}$ and not a fraction of a population one
(Section~\ref{sec:parent}); and $L_0$, the unrecorded part of the load, is carried as an
unknown rather than assumed away, which Section~\ref{sec:excitation} shows is what makes
the identification problem well posed at all.

\paragraph{Why $\alpha$ is not a rate constant.} Elimination proceeds by
conjugating propionate with carnitine, and each molecule eliminated consumes one
molecule of carnitine, which is excreted with the product. In two patients with propionic
acidemia and two with methylmalonic aciduria, urinary excretion of free carnitine was
decreased while excretion of short-chain acylcarnitines, mostly propionylcarnitine, was
increased~\cite{carnitine_excretion_pa}. % C41
Secondary carnitine deficiency follows, and is why supplementation is standard therapy.

\paragraph{Stoichiometric consumption gives a ceiling, not a decline.} Promote carnitine to
a state $c$, let the conjugation reaction draw on it, and replenish the pool at a constant
rate $R$. If conjugation is the only route out of the pool, $\dot c = 0$ fixes the
conjugation flux at $R$ for every $m$: the alternative channel acquires a ceiling and Model
B is recovered. Two things are missing, a reason for the flux to vanish as the metabolite
does, and a reason for it to decline.

\paragraph{One mechanism supplies both.} Carnitine is recovered from the glomerular filtrate
by a renal transporter that acylcarnitines share. In a carnitine transporter solubilised
from renal apical brush-border membrane and reconstituted into liposomes, substrate
analogues inhibited transport, the most effective being acylcarnitines, and besides
carnitine only acylcarnitines and betaine were efficiently
translocated~\cite{octn2_renal_transport}. % C42
Propionylcarnitine rises with the metabolite, so the fraction of filtered carnitine
recovered falls as the metabolite accumulates. The free pool therefore has a second exit,
and that exit grows with $m$.

\paragraph{The model.} Writing $c(t)$ for available carnitine,
\begin{equation}
\begin{aligned}
\dot m &= s(t) \;-\; \theta\, g(m) \;-\; k\, c\, \frac{m}{K_c + m}, \\[2pt]
\dot c &= R_0 \;-\; k\, c\, \frac{m}{K_c + m} \;-\; \delta(m)\, c,
\qquad \delta(m) = \delta_0\left(1 + \frac{m}{M_i}\right),
\end{aligned}
\label{eq:C}
\end{equation}
with $K_c$ the half-saturation of the conjugation step, $\delta_0$ the basal rate of
non-conjugative loss, and $M_i$ the metabolite level at which competition doubles it.
Replenishment $R_0$ is constant. The conjugation term is the one the two equations share,
entering both with the same coefficient, which is the one-to-one stoichiometry: a molecule
of propionate leaves only with a molecule of carnitine. The two equations are otherwise not
symmetric, because carnitine has a second exit that propionate does not. The enzymatic term
$\theta g(m)$ is unchanged, with $\theta$ and $K_m$ carrying the meanings of Section~3.

\paragraph{The reduction.} Solving
$\dot c = 0$ for $c$ and substituting gives a quasi-steady alternative flux
\begin{equation}
\alpha(m) \;=\; R_0\,\frac{u(m)}{u(m) + \rho(m)},
\qquad u(m) = \frac{m}{K_c+m}, \quad \rho(m) = \rho_0\left(1 + \frac{m}{M_i}\right),
\label{eq:alpha}
\end{equation}
with $\rho_0 = \delta_0/k$. It vanishes at $m = 0$ because $u$ does, rises while the
conjugation step fills, and falls once competition outpaces it. The reduction preserves the
constants: $R_0$ and $K_c$ enter \eqref{eq:alpha} exactly as they enter \eqref{eq:C}, with
no rescaling, and $k$ cancels. Everything from Section~\ref{sec:intrinsic} onwards is
carried out on $C(m;\theta) = \theta g(m) + \alpha(m)$.

This paragraph previously opened by asserting that the carnitine state is fast relative to
the metabolite. That assertion is withdrawn, and Section~\ref{sec:reduction} of the paper gives
the justification that replaces it: $\alpha$ is defined by solving $\dot c = 0$, and
$\dot c = 0$ holds exactly at any fixed point, so at equilibrium the reduced and two-state
conditions are one equation and the reduction is exact there whatever the timescales do.
Away from equilibrium the premise is needed and does not hold, carnitine being the slower
state both in the published human kinetics and in this model's own $\tau_c/\tau_m$. What
that costs, and why $k$ cancelling from \eqref{eq:alpha} does not make $k$ a free scale, is
in Table~\ref{tab:paramsfull} and in Section~\ref{sec:reduction} of the paper.

\paragraph{The constants.} Every one of them is in Table~\ref{tab:paramsfull} with its
value and its source, $M_i = 300$ among them. A constant a figure depends on and the text
never states leaves that figure unreproducible, which is why the register is generated from
the model code rather than typed beside it.

One constraint fixes $R_0$ outright. It is set so the alternative route alone peaks at $500$
against a baseline load of $600$, so that a patient with $\theta = 0$ cannot be maintained,
which the severity of null genotypes requires. $R_0$ is therefore \emph{not} a free
constant, and $\alpha$ is exactly proportional to it, so moving $K_c$, $M_i$ or $\rho_0$
moves the peak and $R_0$ must be re-derived. A sensitivity analysis that holds $R_0$ while
moving one of the other three is sweeping parameter sets that violate the model's own
calibration, and reports the sensitivity to that calibration as though it were sensitivity
to the shape constant. Remark~\ref{rem:shape} keeps the two apart.

$K_c$, $M_i$ and $\rho_0$ have no measurement behind any of them. $K_c$ is set equal to
$K_m$, and there is no reason the half-saturation of a conjugation reaction should match that
of an enzyme; $M_i$ is the level at which acylcarnitine competition doubles the
non-conjugative loss, and nothing measures it; $\rho_0 = 0.5$ sits inside a window the next
paragraph leaves open. Remark~\ref{rem:shape} reports what all three cost.

\paragraph{Carnitine excretion, predicted and measured.}
Absolute urinary free carnitine is the filtered load times the fraction not recovered. In
\eqref{eq:C} the filtered load is proportional to $c$ and the loss flux is $\delta(m)\,c$,
so the fraction rises as $1 + m/M_i$ while the load falls with the plasma level. Free
carnitine of $5.7$~$\mu$M against a normal $37 \pm 8$~\cite{carnitine_pa_child} % C25
makes the load $6.5$ times lower, and at the attractor the fraction rises only $1.3$-fold,
so the model predicts absolute excretion at about a fifth of normal. That is the direction
measured~\cite{carnitine_excretion_pa}. % C41
The prediction fails once $1 + m/M_i$ exceeds $37/5.7$, that is above
$m \approx 5.5\,M_i = 1650$, some $1.6$ times the separatrix. This paragraph is the
arithmetic $M_i$ had to be recovered from, and it is now reproducible from
Table~\ref{tab:paramsfull} instead.

The model satisfies the constraint; it does not derive it. What the constraint excludes is
$\rho_0$ below $0.20$, where the separatrix moves far enough out to reverse the sign. Holding
$R_0$ while moving $\rho_0$ gives $0.35$ instead, but that sweep breaks the calibration
above and is the same conflation Remark~\ref{rem:shape} separates. The constraint does not bound $M_i$ at
all, in either direction, over any range worth sweeping.

\begin{figure}[!htbp]
\centering
\resizebox{0.90\linewidth}{!}{%
\begin{tikzpicture}
\node[mpool] (m) at (0,0) {$m$\\[1pt]\scriptsize\mdseries accumulated\\metabolite};
\node[mlabel] (in) at (-3.6,0) {$s(t)$};
\draw[mar] (in) -- (m);
\node[msmall] at (-2.4,0.42) {load};

\node[mpool] (c) at (5.6,0) {$c$\\[1pt]\scriptsize\mdseries available\\carnitine};
\node[mlabel] (rep) at (9.6,0) {$R_0$};
\draw[mar] (rep) -- (c);
\node[msmall] at (8.0,0.42) {replenishment};

\node[menzy] (enz) at (-2.3,-2.7) {$\theta\, g(m)$\\{\scriptsize saturating}};
\channelglyph{(-3.55,-4.15)}{green!45!black}{s}
\draw[mar] (m.south) -- ++(0,-0.85) -| (enz.north);
\node[msmall] (o1) at (-2.3,-5.0) {eliminated};
\draw[mar] (enz) -- (o1);

\node[malt] (kcm) at (2.5,-2.7) {$k\, c\, m/(K_c+m)$\\[1pt]\scriptsize\mdseries consumes one $c$\\per $m$ removed};
\channelglyph{(1.15,-4.15)}{blue!55!black}{n}
\draw[mar, blue!55!black] (m.south east).. controls +(0.7,-0.9).. (kcm.north west);
\draw[mar, blue!55!black] (c.south west).. controls +(-0.7,-0.9).. (kcm.north east);
\node[msmall] (o2) at (2.5,-5.1) {eliminated together,\\as propionylcarnitine};
\draw[mar] (kcm) -- (o2);

\node[malt] (dc) at (6.9,-2.7) {$\delta(m)\, c$\\[1pt]\scriptsize\mdseries not recovered\\from the filtrate};
\draw[mar, blue!55!black] (c.south east).. controls +(0.6,-0.9).. (dc.north);
\node[msmall] (o3) at (6.9,-5.1) {lost in urine};
\draw[mar] (dc) -- (o3);

\draw[red!60!black, thick, -{Bar[width=2.4mm]}]
      (m.north) -- ++(0,1.45) coordinate (fbtop) -| (10.35,-2.7) -- (dc.east);
\node[msmall, red!60!black, fill=white, inner sep=3pt]
     at ([shift={(5.0,0)}]fbtop)
     {rising $m$ means more acylcarnitine competing for the renal\\
      transporter, so less filtered carnitine is recovered};
\end{tikzpicture}}
\caption{The dynamic model \eqref{eq:twostate} as a block diagram. The pool $m$, the load $s(t)$
and the enzymatic channel are placed and coloured as in Figure~\ref{fig:blockA} of the
paper, and
everything else differs. Carnitine is a second state $c$, and the conjugation term is drawn
as a single node draining both pools, because the one-to-one stoichiometry is the
mechanism: each molecule of propionate that leaves takes a molecule of carnitine with it.
Replenishment $R_0$ is constant. The exit on the right is the one carnitine has and
propionate does not, and the blunt-ended arrow is the whole content of the model:
acylcarnitine rising with $m$ competes for the renal transporter, so less filtered carnitine
is recovered and that exit widens. Having a second exit at all is what makes the flux vanish
with the metabolite; the exit widening is what makes it decline rather than saturate. The
glyph under each channel is its shape.}
\label{fig:blockC}
\end{figure}

Figure~\ref{fig:blockC} draws the two-state structure and the widening carnitine exit.

\subsubsection{Consequence: positive feedback}

The loop is
\[
m \uparrow \;\Longrightarrow\; \text{acylcarnitine} \uparrow \;\Longrightarrow\;
\text{less carnitine recovered} \;\Longrightarrow\;
c \downarrow \;\Longrightarrow\; \text{clearance} \downarrow \;\Longrightarrow\; m \uparrow.
\]

Consumption by conjugation is not part of that loop: on its own it draws the pool down and
the replenishment holds it up, which is a ceiling. What closes the loop is the widening exit,
and once it outpaces the conjugation step effective clearance \emph{decreases} with $m$
rather than merely saturating. The clearance curve becomes non-monotone, and the system admits
two equilibria at moderate load: a stable one at low $m$, and an unstable one at
higher $m$. As load rises the two approach, collide and annihilate, which is a genuine
saddle-node.

\subsubsection{Why this would matter clinically}

The unstable equilibrium is a \emph{separatrix}. Beyond it, $\dot m > 0$ even at baseline
load, so the metabolite continues to rise after the precipitating illness has resolved.
This is a statement about the model's state, not about the measured marker, and
Remark~\ref{rem:scope} below says why the two must not be conflated here.
Removing the trigger returns $s$ to its baseline value, and at that value the patient is
still on the wrong side of the separatrix. Recovery therefore requires an active intervention that moves one of the two: the load
driven \emph{below} baseline, which raises the separatrix above the patient, or the
metabolite lowered directly below the separatrix where it stands. These are intravenous
glucose with protein withdrawal, and haemofiltration.

\begin{remark}[Why the question is phrased this way]
The obvious phrasing, ``why does a crisis not resolve spontaneously'', is not what the
model says, and Figure~\ref{fig:C2} contradicts it: lowering the load \emph{does} recover
the patient, because it raises the separatrix. What the model says is that recovery has to be
bought by an active intervention. Removal of the trigger
restores the vector field that the patient has already escaped, so the intervention has to
overshoot, and the amount of overshoot is set by how far above the separatrix the patient
sits. Neither reduced structure predicts this: in both, returning $s$ to baseline
returns $m$ to baseline.
\end{remark}

\subsubsection{What is and is not new here}
\label{sec:priorart}

Two parts of \eqref{eq:parent} have prior art in this disease and one does not, and the
distinction is worth making before anything is proved, because it says what this paper is
claiming.

\paragraph{Therapy as a time-varying input is not new.} Both published pharmacokinetic and
pharmacodynamic models of mRNA-3927 are indirect-response models in which the dose drives a
time-varying quantity rather than a fixed one. The translational model puts the indirect
response on the enzyme and makes the markers a direct saturating function of
it~\cite{pkpd_translational}, % C46
and the population model that selected the biweekly dose treats each marker as a turnover
state whose production the therapy inhibits and whose loss is first
order~\cite{dose_selection_poster}. Writing $\theta + \Delta(t)$ is the same move. What
\eqref{eq:parent} adds is that $\theta$ is the \emph{patient's own} residual capacity and
therefore the object to be identified, rather than a population parameter with the
between-patient variation carried by random effects.

\paragraph{A longitudinal reformulation is not new either.} The natural history study
already fits a hazard of the shape Section~\ref{sec:hazard} proposes: a generalised additive
\emph{mixed} model with a smooth term in days since the previous event and a
participant-level random effect~\cite{natural_history}. % C28
What is proposed there is a substitution \emph{within} that machinery (a mechanistic
$\theta$ in place of a measured biomarker as the covariate carrying the risk) and not
new statistical apparatus.

\paragraph{What is new is the question asked of the input.} Neither prior model asks whether
the dosing schedule \emph{excites} the parameter it is estimating. Both are fitted to a
cohort in which the schedule is fixed by protocol and the parameters are population
quantities, which is a well-posed estimation problem with or without excitation. This paper
asks the individual question, and Section~\ref{sec:excitation} is the answer. A fixed
schedule supplies almost no excitation of the direction separating capacity from
half-saturation \emph{when the schedule is not therapeutic}. A fixed schedule of
\emph{therapeutic} doses excites that direction directly: the sensitivity ratio
$-(1 + \Delta(t)/\theta)\,\varphi(m)$ is modulated by $\Delta(t)$, and the dose train
supplies forty-two times what the marker band does on the same record
(Section~\ref{sec:therapychannel}). What supplies almost none is a record without one, and
the non-accumulation result holds there: Section~\ref{sec:whoto} reports that the
information in the weak direction grows with the number of draws but never arrives, because
the load-record prior holds that direction at every $n$ a study could reach. The one way
back to the original sentence is a sampling schedule aliased against the dose interval,
which is a design error rather than a property of fixed schedules
(Remark~\ref{rem:aliasing}).

\subsubsection{The reduction hierarchy, verified}
\label{sec:hierarchy}

Three objects are in play and they are not three models. They are one model read at three
resolutions, and the relations between them were checked symbolically and numerically
rather than asserted, because one of the three does not hold.

\begin{description}[leftmargin=1.6em, itemsep=3pt]
\item[The dynamic model.] Equation \eqref{eq:parent} with the two-state carnitine system
\eqref{eq:C} behind it. Verified: solving $\dot c = 0$ and substituting returns
\eqref{eq:alpha} exactly, symbolically, with $k$ entering only through $\rho_0 = \delta_0/k$.
\item[The equilibrium reduction.] Setting $\dot m = 0$ gives $L = \theta g(m) + \alpha(m)$,
whose solutions are the attractor and the separatrix, and whose fold defines
$\thc^{\,\mathrm{int}}$. Exact, and by identity rather than by verification: integrating the
two-state system from below the separatrix settles on the reduced attractor to fourteen
digits, but that agreement is $\dot c = 0$ holding at a fixed point and cannot fail at any
parameter value, so it is not the check it was previously listed as. The peak of $C$ at
$\theta = \thc^{\,\mathrm{int}}$ equals the load to ten digits, which is a genuine check
because both sides are computed from $C$ alone. The dynamical half of the old entry, that
integrating from above the separatrix diverges, holds on the slow manifold and fails off it:
with $c$ left at its pre-perturbation value the basin runs $10\%$ to $18\%$ past the reduced
separatrix at the sampled corners, and up to $50\%$ at the extremes of the shape constants.
\item[The local identification model.] Section~\ref{sec:local}, over the band a managed
patient visits. Not verified, and it does not hold as a reduction:
Remark~\ref{rem:not_a_reduction} says what does.
\end{description}

\begin{remark}[The local model is not a reduction of this one, and that is why it is the
favourable case]
\label{rem:not_a_reduction}
It would be convenient if Section~\ref{sec:local} were \eqref{eq:parent} linearised about
the operating point. It is not, and the check is one line. The tangent to
$C = \theta g + \alpha$ at $\bar m$ has intercept
\begin{equation}
a \;=\; \theta\, g(\bar m)^2 \;+\; \bigl[\alpha(\bar m) - \bar m\,\alpha'(\bar m)\bigr],
\label{eq:parent_intercept}
\end{equation}
so the local model's identity $a = \theta g(\bar m)^2$ holds \emph{only} when the second
bracket vanishes, that is only when the alternative channel is a straight line through the
origin. At the illustrative constants the bracket is $374$ against the enzymatic term's
$61$: the alternative channel owns six-sevenths of the intercept, which is the term
the residual capacity was supposed to be carried by.

What the local model is, then, is \eqref{eq:parent} with $\alpha$ replaced by $k_e m$. That
replacement is not an approximation of this model; it is a different hypothesis about one
term. It is retained because it is the \emph{most favourable} case for identification, for a
reason now sharper than the one this paper used to give. A curved $\alpha$ contributes to
both coefficients of an affine fit, so given $\alpha$ the pair $(a,b)$ determines
$(\theta, K_m)$ outright, verified to machine precision, and the fibre of
Remark~\ref{rem:fibre} does not exist. A straight $\alpha$ contributes to neither, and the
fibre appears. The degeneracy the paper is about is a property of the most favourable
structure, so it is not an artefact of the awkward one.

The proviso carries weight in the other direction. ``Given $\alpha$'' means given four
constants, and Remark~\ref{rem:shape} prices them. What the record does supply is a
two-sided bound: fitting $(a,b)$ excludes all but roughly $K_c \in [100, 300]$, and inside
that window the recovered $\theta$ moves by a factor of $1.41$.
\end{remark}

\begin{remark}[The hypothesis with a ceiling instead of a fold]
\label{rem:ceiling}
A third structure is a single hypothesis about $\alpha$ rather than a separate model. Let
the alternative
channel saturate rather than decline: $\alpha(m) = k_e h(m)$ with $h(m) = m/(K_e+m)$. Then
total clearance is bounded by $\theta + k_e$, loads below that ceiling admit one stable
equilibrium and loads above it admit none.

It is kept for exactly two purposes and neither needs a section.

First, it is the null against which the fold is testable. Under a ceiling, removing
the precipitating illness returns the load to baseline and the marker follows it back; under
a fold, the patient who has crossed the separatrix stays across it and recovery has to be
bought by driving the load \emph{below} baseline. That is the discriminant of
Section~\ref{sec:discriminant}, and it has no content without the null.

Second, it prices biological fidelity in the identification, and it does so twice
over. Once both channels saturate they share a functional form and are separated only by
bending in different places, by $K_e \neq K_m$, and at a ten per cent excursion the
condition number rises from $1.1\times10^{5}$ to $4.8\times10^{7}$, a factor of about
$450$. That is the \emph{local} price and it is the one this remark originally quoted.

There is a second price the condition number cannot show, because it is not a matter of
conditioning at all. By Proposition~\ref{prop:exchange} two channels of shared functional
form are exchangeable, so the ceiling hypothesis carries an exact two-point degeneracy in which
$(\theta, K_m)$ and $(k_e, K_e)$ swap: the Jacobian stays full rank, which is why a finite
condition number is the right local description, and a second isolated parameter vector
nevertheless reproduces the data exactly. A selective, \emph{time-varying} input removes it
and a constant one does not. The paper's own $\alpha$ is not of this form: \eqref{eq:alpha}
is non-monotone and $\theta g(m)$ is not, so they are not two members of one family and no
exchange is available, but the ceiling hypothesis is, and a reader who preferred it would
inherit the degeneracy along with it.

No equilibrium collides with another here; a single equilibrium escapes to infinity as the
load approaches the ceiling. A genuine saddle-node needs a non-monotone $\alpha$, which is
what \eqref{eq:alpha} supplies.
\end{remark}

\subsubsection{An intrinsic definition of $\thc$}
\label{sec:intrinsic}

The local reduction defines the critical capacity by \eqref{eq:thcrit}, and the hypothesis
with a ceiling (Remark~\ref{rem:ceiling}) by its own counterpart. Both import two quantities
that are never measured: $\mtol$, a clinical convention about what metabolite level is tolerable, and
$\spk$, the amplitude of an illness. The convention is the more troubling of the two,
because it is not a fact about the patient at all. The equilibrium reduction does not need it.

Write $C(m;\theta) = \theta g(m) + \alpha(m)$ for the total clearance, with $\alpha$ the
carnitine-mediated flux \eqref{eq:alpha}. Since $\alpha$ rises, peaks and then decays,
$C(\cdot\;\theta)$ has an interior maximum for every $\theta$, and that maximum increases
with $\theta$.

\begin{proposition}[Intrinsic critical capacity]
\label{prop:intrinsic}
Fix a load $s$. Define
\begin{equation}
\thc^{\,\mathrm{int}}(s) \:=\; \text{the unique } \theta \text{ such that }
\max_{m > 0}\, C(m;\theta) \;=\; s.
\label{eq:thcrit_int}
\end{equation}
Then for $\theta > \thc^{\,\mathrm{int}}(s)$ the system has exactly two equilibria at load
$s$, a stable attractor and an unstable separatrix, and for
$\theta < \thc^{\,\mathrm{int}}(s)$ it has none, so the metabolite diverges from every
initial condition. At $\theta = \thc^{\,\mathrm{int}}(s)$ the two equilibria merge in a
saddle-node.
\end{proposition}

\begin{proof}[Argument]
Equilibria solve $C(m;\theta) = s$. The map $\theta \mapsto \max_m C(m;\theta)$ is
continuous and strictly increasing, since $C$ is strictly increasing in $\theta$ pointwise
and the maximum is attained in the interior, so \eqref{eq:thcrit_int} has a unique
solution. If the maximum exceeds $s$ the horizontal line $s$ cuts the graph of $C$ twice,
once on the rising branch, where $C$ is increasing and the equilibrium is stable, and once
on the falling branch, where it is unstable. If the maximum is below $s$ there is no
intersection and $\dot m = s - C > 0$ everywhere. At equality the two intersections
coincide at the maximum, where $C' = 0$, which is the saddle-node condition.
\end{proof}

Two consequences, independent of each other:

\begin{itemize}[leftmargin=*, itemsep=3pt]
\item \label{item:wellposed}\textbf{$\mtol$ has disappeared.} No tolerability convention enters
\eqref{eq:thcrit_int}. The threshold is the value of $\theta$ at which the model's own
equilibrium structure changes, and it would be the same threshold under any convention
about which metabolite levels a clinician is willing to accept. Under the reduced structures no such
definition is available, because neither has a second equilibrium to collide with: in A the
equilibrium is never lost, and in B it escapes to infinity without annihilating. This is
the sense in which only the equilibrium reduction makes $\thc$ well posed, and
it is the reason the model is in the paper at all. Its biological fidelity is the usual
reason given and the weaker one.
\item \textbf{The margin is the clinically decisive object.} Above the threshold, what
matters is how far apart the two equilibria are. The
distance from the attractor to the separatrix is how far the patient can be pushed before
the crisis becomes self-sustaining, and it closes continuously to zero at
$\thc^{\,\mathrm{int}}$. That distance is computed and plotted in
Figure~\ref{fig:C5}. It has no counterpart under either reduced structure.
\end{itemize}

\begin{corollary}[When the threshold exists]
\label{cor:fold}
$\thc^{\,\mathrm{int}}(s) > 0$ if and only if $\max_m \alpha(m) < s$. The fold exists only
while the alternative route alone cannot carry the load. Where it can, a patient with no
residual enzyme is maintainable, the fold sits at negative $\theta$, and the threshold is
vacuous.
\end{corollary}

This is what fixes $R_0$. At the illustrative constants the alternative route peaks at $500$
against a baseline load of $600$, giving $\thc^{\,\mathrm{int}} = 151$; raising $R_0$ far
enough to lift that peak past $600$ would make null genotypes maintainable, which the
clinical picture excludes.

\paragraph{The excursion an illness has to make.} The fold exists to say why recovery needs
more than removal of the trigger, and that requires an illness to be able to strand the
patient in the first place. A sustained load below the fold cannot. At any
$s < s_{\mathrm{fold}}$ the system settles on an attractor lying below the separatrix at
baseline load, so returning $s$ to baseline returns the patient with it. Crossing requires
the load to exceed $s_{\mathrm{fold}}$, the largest load the patient can balance at all,
which at $\theta = 300$ is $706$ against a baseline of $600$. An excursion of
$18\%$ in the load is the smallest that strands a patient, and the illnesses drawn in
Figures~\ref{fig:C1} and \ref{fig:C2} are $42\%$ and $50\%$ above baseline, clearing it by
factors of $2.4$ and $2.8$. The margin of Figure~\ref{fig:C5} is measured in metabolite and
is large; the excursion that consumes it is measured in load and is not.

\begin{remark}[What the three unmeasured shape constants cost, and where the cost actually
falls]
\label{rem:shape}
$K_c$, $M_i$ and $\rho_0$ have the same status: chosen, with no measurement behind any of
them. Bounding $K_c$ alone would leave the impression that it is the expensive one. It is
not, and the ordering is not the interesting part of the answer.

\begin{center}
\scriptsize
\begin{tabular}{@{}lllllll@{}}
\toprule
\textbf{constant} & \textbf{range swept} & \textbf{range} & \textbf{$\thc$, $R_0$ re-derived} & \textbf{$\thc$, $R_0$ held} & \textbf{vacuous} & \textbf{margin, $R_0$ re-derived} \\
\midrule
$K_c$ & 25--800 & 32$\times$ & 1.56 & 6.27 & 1 & 3.69 \\
$M_{i}$ & 150--1200 & 8$\times$ & 1.37 & 21.25 & 1 & 6.44 \\
$rho_{0}$ & 0.35--1 & 3$\times$ & 1.02 & 967.47 & 0 & 1.48 \\
\bottomrule
\end{tabular}
% factors are max/min over the swept range; 'vacuous' counts sweep points at which a null-genotype patient becomes maintainable and the threshold ceases to exist
\end{center}

Three things follow, and the second is the one that matters.

First, on a common footing all three are cheap, and $M_i$ is cheaper than $K_c$.
Holding the calibration that fixes $R_0$, an eight-fold range in $M_i$ moves
$\thc^{\,\mathrm{int}}$ by a factor of $1.37$ against $K_c$'s $1.56$ over thirty-two-fold,
and $\rho_0$ across its whole admissible window moves it by $1.6\%$. Sweeping the three
together over those ranges leaves $\thc^{\,\mathrm{int}}$ inside a factor of $1.4$.
The threshold is very nearly a function of the peak of $\alpha$ alone, so its
exposure is concentrated not in three constants nobody has measured but in one clinical
claim: that a patient with no residual enzyme cannot be maintained at baseline load. That
claim can be argued with a clinician. The shape constants cannot.

The column beside it is what happens when $R_0$ is held instead. The apparent sensitivity
inflates by an order of magnitude and the ordering reverses, $\rho_0$ becoming the worst of
the three by a factor of nine hundred, and at one end of each sweep the threshold goes
vacuous. None of that is sensitivity to the shape constant; it is the peak of $\alpha$
crossing the baseline load, which the calibration exists to prevent.

Second, the margin is not protected the way the threshold is. The separatrix sits
where $\alpha$ has decayed, and $\alpha$ decays on the scale of $M_i$, so the margin
inherits $M_i$ almost linearly where $\thc^{\,\mathrm{int}}$ does not: over the same
eight-fold range it runs from $509$ to $3275$. Figure~\ref{fig:C5} now draws that
range as a band rather than quoting the nominal $954$ as a curve, so the figure carries its
own uncertainty instead of delegating it to this remark. What survives the sweep is the
\emph{sign}: the margin is positive at every $\theta$ above the threshold and at every
$M_i$ in the range. What does not survive is the value, and the quantity
Section~\ref{sec:intrinsic} calls ``the clinically decisive object'' is the one this model
pins down least well.

Third, the well-posedness claim survives and the quantitative claim does not.
Proposition~\ref{prop:intrinsic} is a well-posedness statement and no shape constant touches it.
What the constants move is where the threshold sits and how much room is left above it.
\end{remark}

\begin{remark}[The limit of the construction]
\label{rem:intrinsic_limit}
The construction eliminates $\mtol$. It does not eliminate $\spk$. Equation
\eqref{eq:thcrit_int} takes the load as an argument, and the clinically operative question
is not whether a patient is maintainable at baseline load but whether an intercurrent
illness pushes them across the separatrix. That depends on how far and for how long the
load is raised, which is $\spk$ and the episode duration, and neither is measured. What has
changed is where the unmeasured quantity sits: it is no longer inside the definition of the
threshold, where it corrupted the threshold itself, but in the separate question of whether
a given margin is large enough. The remaining step is to stop treating episodes as a single
worst case and give their amplitudes a distribution estimated from the records, which turns
the deterministic crossing into a hazard. That is the reformulation set out in Section~7,
and only it removes $\spk$.
\end{remark}

\subsubsection{What would have to be established}
\label{sec:notworking}

Four things: that replenishment is impaired \emph{enough}
to turn clearance downward, which is a question about within-patient free-carnitine
trajectories and not about group means; that replacing $k_e$ with a channel carrying four
constants is a net gain, which Remark~\ref{rem:shape} prices and
Remark~\ref{rem:not_a_reduction} bounds; that a competition documented in a reconstituted
rat renal transporter carries to patients; and that adding one state to capture half a
mechanism is better than not adding it.

\subsubsection{Testable discriminant, and the observable it has to be stated in}
\label{sec:discriminant}

The fold and the ceiling differ in whether removing the trigger suffices for recovery, and
that difference is a statement about $m$ which \emph{cannot be taken to the records in that
form}: propionylcarnitine, propionylglycine and propionic acid are reported essentially
unaltered through an acute decompensation~\cite{amd_phenotyping}. % C24
The single marker does not move through the excursion the discriminant is about.

Free carnitine does move, and it is measured routinely, so the discriminant is restated in
it: the median falls from $38.9$ to $22.0$~$\mu$M between samples drawn outside a
decompensation and during one, $P = 0.0370$~\cite{amd_phenotyping}. % C38
The phenotyping study's own discussion independently advises measuring free carnitine during
a decompensation~\cite{amd_phenotyping}, % C39
so a model built to represent depletion and a recommendation written without reference to
any model arrive at the same analyte.

Two qualifications: the comparison pools
samples across the cohort rather than within patients, and a fall is consistent with the
fold rather than exclusive to it.

\begin{remark}[Scope]
\label{rem:scope}
Sections~\ref{sec:local} onwards concern the \emph{stable} patient, where the marker and the
state can be argued to move together. The behaviour through a decompensation is exactly
where the measured proxy stops tracking the state, so the separatrix is a claim about the
dynamics that the current observable cannot confirm or refute directly, which is also why
Section~\ref{sec:marker_orig} fixes the marker for the stable patient and not for the crisis.
\end{remark}

\vspace{0.5em}
\hrule
\vspace{1em}

%==========================================================================
